\begin{abstract}
We introduce bindings that enable the convenient, efficient, and
reliable use of software modules of \cgal{} (Computational Geometry
Algorithm Library), which are written in C++, from within code written
in Python. There are different tools that facilitate the creation of
such bindings. We present a short study that compares three main
tools, which leads to the tool of choice. The implementation of
algorithms and data structures in computational geometry presents
tremendous difficulties, such as obtaining stable software despite the
use of (inexact) floating point arithmetic, found in standard
hardware, and meticulous handling of all degenerate cases, which
typically are in abundance. The code of \cgal{} extensively uses
function and class templates in order to handle these difficulties,
which implies that the programmer has to make many choices that are
resolved during compile time (of the C++ modules). While bindings
take effect at run time (of the Python code), the type of the C++
objects that are bound, must be known when the bindings are generated,
that is, when they are compiled. The types of the bound objects are
instances (instantiated types) of C++ function and class
templates. The number of object types that can potentially be bound,
in implementation of generic computational-geometry algorithms, is
enormous; thus, the generation of the bindings for all these types in
advance is practically impossible. For example, the programmer needs
to choose among a dozen types of curves (e.g., line segments, circular
arcs, geodesic arcs on a sphere, or polycurves of any curve type) to
yield a desired arrangement type; often there are several choices to
make, resulting in a prohibitively large number of combinations. We
present a system that dynamically generates bindings for desired
object types according to user prescriptions, which enables the
convenient use of any subset of bound object types concurrently. After
many years, in which the usage of these packages was restricted to C++
experts, the introduction of the bindings made them easily accessible
to newcomers and practitioners in non-computing fields, as we report
and discuss in the paper.
\end{abstract}
